\section{Open Questions, Ambiguities, and Fragile Claims}

\subsection{Concepts that remain fuzzy}

Several important concepts are useful but not fully specified in the source material:

\begin{itemize}[leftmargin=1.5em]
\item what exactly separates a ``real'' density from a fake one beyond trader intuition;
\item how many tests are ``too many'' before a bounce setup should be downgraded into breakout risk;
\item how to quantify a ``clean'' tape versus a noisy one;
\item how to tell when a robot is meaningful enough to trade with;
\item how much of Smart Money language is explanatory metaphor versus testable mechanism.
\end{itemize}

\subsection{Potential contradictions}

The material contains productive tension rather than outright incoherence:

\begin{itemize}[leftmargin=1.5em]
\item densities are treated as important, yet repeatedly described as insufficient on their own;
\item breakouts are described as dangerous for beginners, yet also necessary for advanced opportunity capture;
\item historical levels and POC zones matter, yet the same sources insist that today's context can completely override yesterday's map;
\item robots are tradable when obvious, but their increasing sophistication makes simple pattern-copying unreliable.
\end{itemize}

These are not flaws so much as reminders that the real object of study is \emph{interaction}, not static pattern shape.

\subsection{Claims most in need of empirical testing}

The following claims deserve explicit measurement rather than trust:

\begin{enumerate}[leftmargin=1.5em]
\item round-number densities are materially more reliable than irregular-price densities;
\item replenishing visible queues predict reversal more strongly than merely large static queues;
\item early-session behavior is enough to build same-day setup efficacy filters;
\item robot signatures visible to humans can be captured robustly enough to improve signals;
\item cluster-backed levels outperform visually similar levels without historical traded-volume confirmation.
\end{enumerate}

\subsection{Claims that may be outdated or venue-specific}

Some of the source logic may depend on the era or venue:

\begin{itemize}[leftmargin=1.5em]
\item old MOEX-style order-book behavior may not transfer cleanly to current crypto;
\item current exchange matching, hidden-order behavior, and retail participation may differ from the environments in which the discretionary heuristics matured;
\item planka mechanics, prop controls, and specific terminal affordances are venue-specific and should not be generalized automatically;
\item some ``robot'' patterns may disappear as execution algorithms become less visually naive.
\end{itemize}

\subsection{Source gaps}

The missing volatility-harvesting and psychology transcripts create two blind spots:

\begin{itemize}[leftmargin=1.5em]
\item the exact intended formulation of ``volatility harvesting'' is only partially inferable from neighboring materials;
\item the explicit teaching on trader psychology is absent, so the psychology section here relies mostly on the risk-management book and scattered comments from other lessons.
\end{itemize}

These gaps should eventually be filled with the original source files or new repository notes.
